\documentclass[11pt]{article}
\usepackage[a4paper,margin=23mm]{geometry}
\usepackage[T1]{fontenc}
\usepackage{lmodern,amsmath,amssymb,amsthm,mathtools,microtype,booktabs}
\usepackage[colorlinks=true,linkcolor=blue,citecolor=blue,urlcolor=blue]{hyperref}
\newcommand{\EE}{\mathbb E}
\newcommand{\Om}{\Omega_O}
\newcommand{\ph}{\varphi}
\newcommand{\bn}[2]{\binom{#1}{#2}}
\newcommand{\fall}[2]{(#1)_{\underline{#2}}}
\newtheorem{theorem}{Theorem}
\title{All fixed joint words above a Slater state\\
\large Local CAR polynomials, the same-disorder limit, and finite-time evolution}
\author{Alan Day \and Chengheng\thanks{AI research assistant (OpenAI).}}
\date{7 November 2026\quad Local research draft}
\begin{document}\maketitle
\begin{abstract}
We extend the same-disorder Slater calculation from fourth moments
to every fixed joint word of the Hamiltonian and its selected
particle--hole block. An exact single-site polynomial retains CAR
signs, including negative finite-size six-point contributions.
A coefficient estimate controls all multiple-overlap configurations
after normalization. Only quadratic site clusters survive, while
backward selected edges vanish. The resulting positive word law
constructs the joint cyclic representation. Uniform finite-size
row/column covariance bounds then control the complete time series,
giving annealed finite-time correlators and a self-adjoint limiting
Hamiltonian. This is a model-level joint dynamics result; the
stationary physical state and the de Sitter bulk dictionary remain
separate targets.
\end{abstract}

\section{The model and what is being extended}
Use $N$ complex fermions with standard CAR, $m=N/2$, $a=N-m=m$,
and a fixed occupied set $O$, empty complement $E$, and
$\Om=|O\rangle$. With increasing-order creators and adjoint
annihilators, set
\begin{equation}
 H_N=\sqrt{w_N}\sum_{|A|=|B|=p}J_{AB}C_A^\dagger C_B,
 \quad w_N=\frac{4^pJ_s^2}{\lambda_N}\left(\frac{p!}{N^p}\right)^2,
 \quad\lambda_N=\frac{p^2}{N}.
\end{equation}
The centered Hermitian Gaussian matrix obeys
$\EE J_{AB}J_{CD}=\delta_{AD}\delta_{BC}$; overlapping $A,B$
are included. Write $\ph_N(X)=\EE\langle\Om,X\Om\rangle$.
All operators below use this same $J$. Take $p\to\infty$,
$\lambda_N\to\lambda>0$, fixed $J_s>0$, $\hbar=1$ and string
units for $H_N$. The microscopic state is not assumed stationary.

\newpage
The normalized operators and exact counts are
\begin{align}
 L_N&=\bn mp\bn ap,&D_N&=\bn mp\bn{a+p}p,\\
 h_N&=H_N/\sqrt{w_ND_N},&
 M_N^\dagger&=L_N^{-1/2}
       \sum_{P\subset E,\ Q\subset O}J_{PQ}C_P^\dagger C_Q.
\end{align}
All subsets have size $p$, and $M_N$ is the adjoint. Define
\begin{equation}
 q=e^{-4\lambda},\qquad s=\sqrt q,\qquad\gamma=q^{1/4}.
\end{equation}
The preceding stage evaluated all words through degree four.
It found $\ph(hM^\dagger)=\gamma$ and a seven-dimensional
degree-at-most-two quotient. The task here is the full joint law,
not another derivation of the pure-$M$ dressed oscillator.

The literature inputs are the exact oriented CAR sign organization
and pair-intersection method of complex DSSYK
\cite[Sec.~3.2--3.3]{BNR}, and the state-adapted perspective of
Rajgadia--Xu \cite{RX}. Their separate matter disorder is not
substituted for our same-$J$ block. Our coefficient proof below
keeps the fixed Slater occupation instead of a trace and bounds
the discarded clusters after the required normalization.

\section{The limiting word rule}
Write a word in $h,M^\dagger,M$ from left to right. For each
Gaussian pairing $\pi$, call an $hh$ edge $G$. Every other edge
is selected, $S$, and its $M$ endpoint fixes its creation direction.
An edge with equal non-$h$ endpoints is incompatible. A selected
edge is forward when its creation acts first, from the right;
otherwise it is backward. Let $b(\pi)$ count mixed $hM$ edges.
For a pairing with all selected edges forward, let
\begin{equation}\label{nu}
 \nu(\pi)=\#\{\text{crossing edge pairs}\}
 +\#\{\text{nested edge pairs with selected inner edge}\}.
\end{equation}
Disjoint pairs and nestings with general inner edge add zero.

\begin{theorem}\label{words}
For each fixed word $w$ of length $2k$,
\begin{equation}\label{law}
 \lim_N\ph_N(w)=
 \sum_{\substack{\pi\text{ compatible}\
                   \text{all selected edges forward}}}
       \gamma^{b(\pi)}q^{\nu(\pi)}.
\end{equation}
Odd moments vanish. The empty word has value one.
\end{theorem}
The remainder of the proof establishes this for arbitrary fixed
$k$, including signs and the suppression of backward edges.
It does not rely on inferring general factorization from fourth
or sixth moments.

\newpage
\section{An exact single-site polynomial}
Orient each Gaussian edge $e=(i,j)$, $i<j$, with right operator
$X_e=C_{P_e}^\dagger C_{Q_e}$ and left $X_e^\dagger$.
General edges allow arbitrary $p$-subsets $P_e,Q_e$. Forward
selected edges require $P_e\subset E,Q_e\subset O$; backward
ones require $P_e\subset O,Q_e\subset E$.
Introduce $2k$ membership variables $z_{P_e},z_{Q_e}$.

For a site of initial occupation $o\in\{0,1\}$ and a membership
subset $S$ of those $2k$ labels, calculate its one-mode CAR
matrix element by following the slots from right to left.
At a right endpoint annihilate $Q_e$ then create $P_e$; at a
left endpoint annihilate $P_e$ then create $Q_e$. Omit actions
whose labels are absent from $S$. Set the result to zero if a
selected-edge site restriction is violated. Multiply by
\begin{equation}\label{sign}
 (-1)^{\sum_{e<f:\ e,f\text{ cross}}t_e(S)t_f(S)},
 \qquad t_e(S)=\boldsymbol1_{P_e\in S}+\boldsymbol1_{Q_e\in S}.
\end{equation}
Call this coefficient $f_o(S)\in\{-1,0,1\}$, with $f_o(\varnothing)=1$,
and define
\begin{equation}
 F_o(z)=\sum_{S\subset\{P_1,Q_1,\ldots,P_k,Q_k\}}
               f_o(S)\prod_{i\in S}z_i.
\end{equation}
The raw Gaussian contraction is exactly
\begin{equation}\label{exact}
 A_N(\pi)=[\prod_{i=1}^{2k}z_i^p]\ F_0(z)^a F_1(z)^m.
\end{equation}

\subsection*{Why the CAR sign is exact}
Split each number-conserving edge into its two thin, oriented
fermion chords. Grouping events by site produces the sign
$(-1)^{\sum_{u,v\text{ cross}}(p^2-|I_u\cap I_v|)}$.
This is the exact grouping identity in the oriented CAR derivation
\cite[Eqs.~(3.11)--(3.12)]{BNR}; changing the trace to a diagonal
occupation matrix element changes the local factors, not this
permutation sign. The two thin chords within one edge are nested.
Two crossing thick edges give four thin crossings, so their
$4p^2$ contribution is even. The remaining overlap parity is
precisely \eqref{sign}. Reversal phases of the ordered monomials
cancel between the two endpoints. This argument includes sites
belonging to three or more thin labels; they have not been dropped.

Choosing a membership subset at each of the $a$ empty and $m$
occupied sites gives \eqref{exact}. The coefficient enforces size
$p$ for every set. There is no independent-probe replacement and
no positivity assumption on individual finite-size contractions.

\newpage
\section{A normalized coefficient bound for every cluster}
Put $n=a=m=N/2$. Irrespective of edge type, a nonzero singleton
can only be $P_e$ at an empty site or $Q_e$ at an occupied site.
For general and forward selected edges these singletons have
coefficient one. For backward selected edges both are forbidden.
Use the fictitious singleton baseline $L_N^k=\bn np^{2k}$ even
in the backward case.

For each reservoir $o$ and membership $S$ with $|S|\ge2$, let
$n_{o,S}$ be its number of sites. Let
\begin{equation}
 d_i=\sum_{o,S\ni i}n_{o,S},\quad
 A_o=\sum_{|S|\ge2}n_{o,S},\quad
 U_o=kp-\sum_{i:\ r_i=o}d_i+A_o,
\end{equation}
where $r_{P_e}=0,r_{Q_e}=1$. The correct singleton multiplicity
of label $i$ is $p-d_i$. Feasible configurations have $d_i\le p$;
if that singleton is forbidden, feasibility also requires $d_i=p$.
After division by $L_N^k$, each configuration has signed weight
\begin{equation}\label{cluster}
 \frac{\fall n{U_0}\fall n{U_1}}{\fall np^{2k}}
 \prod_i\fall p{d_i}
 \prod_{o,S}\frac{f_o(S)^{n_{o,S}}}{n_{o,S}!}.
\end{equation}
This follows directly from the multinomial expansion of
\eqref{exact}; singleton factorials have been absorbed into
$\fall p{d_i}$.

For fixed $k$, take $kp<n/2$. Since $|f_o(S)|\le1$,
$\fall p d\le p^d$, and the denominator falling factorials
differ from $n^{2kp}$ by a uniformly bounded exponential,
the absolute value of \eqref{cluster} is at most
\begin{equation}\label{bound}
 C_{k,\lambda}\prod_{o,S}
 \frac{\alpha_{S,N}^{n_{o,S}}}{n_{o,S}!},\qquad
 \alpha_{S,N}=n(p/n)^{|S|}.
\end{equation}
Here $C_{k,\lambda}$ can be chosen uniformly along the sequence
with $p^2/N$ in a compact positive interval. In particular,
\begin{equation}
 \sum_{o,|S|\ge2}\alpha_{S,N}=O_k(1),\qquad
 \sum_{o,|S|\ge3}\alpha_{S,N}=O_k(p^3/N^2)=o(1).
\end{equation}
Summing \eqref{bound} over all counts gives an exponential series.
The total absolute weight with at least one higher cluster is
$o(1)$, and the tail with more than $R$ defects is bounded by
$C\sum_{j>R}A^j/j!$ for a fixed finite $A$. This controls signed
cancellations and rare large-overlap configurations after normalization.

\newpage
\subsection*{Forward limit and backward suppression}
For any fixed collection of double-cluster counts, the factorials
in \eqref{cluster} have limit
\begin{equation}
 e^{-2\lambda k(k-1)}
 \prod_{o,|S|=2}\frac{(2\lambda f_o(S))^{n_{o,S}}}{n_{o,S}!}.
\end{equation}
Indeed $U_o=kp+O(1)$, and
$\log\fall n r=r\log n-r(r-1)/(2n)+O(r^3/n^2)$ uniformly
for $r\le kp$; also $\fall p d/p^d\to1$ for fixed $d$.
The bound on defect tails justifies summing these limits:
\begin{equation}\label{forward}
 \frac{A_N(\pi)}{L_N^k}\longrightarrow
 \exp\left[2\lambda\left(
       \sum_{o,|S|=2}f_o(S)-k(k-1)\right)\right].
\end{equation}
For a backward edge, at least one forbidden correct singleton
forces $d_i=p$. Since each defect contains label $i$ at most once,
the number of defect sites is at least $p$. Its total normalized
weight is bounded by $C\sum_{j\ge p}A^j/j!\to0$.
This includes general $H$ edges, for which the old pure-ladder
backward estimate alone would not suffice.

\subsection*{The complete quadratic local table}
For an all-forward graph the sum of double-cluster coefficients
is obtained from the following one-site CAR cases:
\begin{center}
\begin{tabular}{lc}
\toprule
Membership sources & Sum over allowed reservoir and two labels\\
\midrule
Within a general edge: $P_e,Q_e$ & $1$\\
Within a selected edge & $0$\\
Two disjoint edges & $2$\\
Two crossing edges & $0$\\
Nested edges, inner general & $2$\\
Nested edges, inner selected & $0$\\
\bottomrule
\end{tabular}
\end{center}
For disjoint edges the surviving pairs are $P_eP_f$ at empty
sites and $Q_eQ_f$ at occupied sites. For nested edges the two
survivors instead put the inner edge's annihilator on an empty
site or its creator on an occupied site; an inner selected edge
forbids both. Crossings have no alternating two-label local
history. Within a general edge only an initially occupied
site carrying both $P_e,Q_e$ survives.

Writing $k_G$ for the number of general edges, the sum is
$k_G+k(k-1)-2\nu(\pi)$. Equation \eqref{forward} is therefore
$e^{2\lambda k_G}q^{\nu(\pi)}$. A word with $n_h$ Hamiltonian
endpoints has normalization $L_N^k(D_N/L_N)^{n_h/2}$, with
$D_N/L_N\to e^{2\lambda}$ and $n_h=2k_G+b(\pi)$.
The remaining factor is $e^{-\lambda b}=\gamma^b$.
Summing the finitely many pairings proves Theorem~\ref{words}.

\newpage
\section{Uniform finite-size bounds before taking time limits}
Fixed-degree convergence by itself does not justify a real-time
series. Here the actual finite operators have a stronger bound.
Expand all of them in the same independent standard real Gaussian
coordinates $x_\alpha$:
\begin{equation}
 h_N=\sum_\alpha x_\alpha A_\alpha,\qquad
 M_N=\sum_\alpha x_\alpha B_\alpha.
\end{equation}
On the $m$-particle sector,
\begin{equation}\label{rowcol}
 \sum A_\alpha^\dagger A_\alpha=
 \sum A_\alpha A_\alpha^\dagger=I,\qquad
 \sum B_\alpha^\dagger B_\alpha\le I,\quad
 \sum B_\alpha B_\alpha^\dagger\le I.
\end{equation}
The first identity is the exact sector variance $D_N$.
For the selected block the two sums are diagonal on occupation
states. On a state with $x$ particles in $E$ their eigenvalues,
in either order, are
\begin{equation}
 \frac{\bn xp^2}{L_N},\qquad
 \frac{\bn{m-x}p^2}{L_N},
 \qquad 0\le x\le m,
\end{equation}
both at most one. The circular Gaussian real and imaginary
parts have their $1/\sqrt2$ factors; hence no extra factor two
appears. Adjoints obey the same row/column bounds.

\subsection*{Every Wick contraction is a contraction}
For a coefficient family $C_\alpha$, opening a Gaussian edge
uses $v\mapsto\sum_\alpha e_\alpha\otimes C_\alpha v$;
its norm is $\|\sum C_\alpha^\dagger C_\alpha\|^{1/2}$.
Closing it uses $\sum e_\alpha\otimes v_\alpha\mapsto
\sum C_\alpha v_\alpha$, with row norm
$\|\sum C_\alpha C_\alpha^\dagger\|^{1/2}$.
Keep tensor slots for all open Gaussian indices and permute
them unitarily to close the required edge. Following any pairing
from right to left reproduces its matrix product in its original
order. By \eqref{rowcol} its operator norm is at most one.
Thus for every finite $N$ and every word of length $2r$,
\begin{equation}\label{allbound}
 |\ph_N(w)|\le(2r-1)!!.
\end{equation}
This proof allows crossings and the shared disorder. It is not
a scalar-independence approximation. For a word $v$ of length $d$,
\begin{equation}\label{vectorbound}
 \|v\Om\|_{L^2(J)}\le F_d,
 \qquad F_d=\sqrt{(2d-1)!!},\quad F_0=1.
\end{equation}
The bound is uniform in both $N$ and degree. Unlike the cluster
proof, it has no degree-squared exponential prefactor.

\newpage
\section{Joint representation and finite-time convergence}
Define $\ph$ on the free star-polynomial algebra by
Theorem~\ref{words}. It is positive because for every fixed
polynomial $P$,
\begin{equation}
 \ph(P^*P)=\lim_N\ph_N(P_N^*P_N)\ge0.
\end{equation}
Quotient by its null left ideal and complete to obtain the cyclic
GNS space. Left multiplication defines $h=h^*$ and the mutually
adjoint-on-the-core $M,M^\dagger$. This provides an all-degree
constructive representation through its Gram matrices and
multiplication, without assuming a finite closed Hamiltonian.

If $P$ has degree at most $d$ and coefficient one-norm $C_P$,
\eqref{vectorbound} gives
\begin{equation}
 \|h^nP\Omega\|\le C_P F_{n+d},\qquad
 \sum_{n\ge0}\frac{|z|^n}{n!}\|h^nP\Omega\|<\infty
 \quad(z\in\mathbb C).
\end{equation}
Every core vector is entire analytic for $h$. Nelson's analytic
vector criterion makes this symmetric Hamiltonian essentially
self-adjoint; we use the standard theorem as recalled in
\cite[Theorem~4.8]{Neeb}. Its closure generates the unitary
$e^{ith}$. Polynomial insertions on evolved core vectors are
defined by the same graph-norm-convergent series; the moment
bound remains valid with their additional fixed degree.

More concretely, for fixed star-polynomials $P_0,\ldots,P_\ell$,
fixed complex $z_1,\ldots,z_\ell$, and $d=\sum_j\deg P_j$,
the ordered correlator
\begin{equation}\label{time}
 \ph_N\left(P_{0,N}e^{z_1h_N}P_{1,N}\cdots
                       e^{z_\ell h_N}P_{\ell,N}\right)
\end{equation}
converges to the series obtained by replacing every polynomial
word coefficient with \eqref{law}. With
$T=\sum_j|z_j|$ and $C=\prod_j C_{P_j}$, the absolute tail
after total exponential degree $R-1$ is bounded uniformly in $N$ by
\begin{equation}\label{tail}
 C\sum_{n\ge R}\frac{T^n F_{n+d}}{n!}.
\end{equation}
The multinomial theorem combines all time powers into $T^n/n!$.
Their successive majorant ratio is
\begin{equation}
 \rho_n=\frac{T\sqrt{2(n+d)+1}}{n+1}\longrightarrow0.
\end{equation}
It decreases with $n\ge0$. When $\rho_R<1$, \eqref{tail} is
at most $C T^R F_{R+d}/[R!(1-\rho_R)]$.
First take $N\to\infty$ at fixed $R$, then $R\to\infty$;
the bound is uniform on compact complex-time sets. Physical
string time uses $H_N=\sqrt{w_ND_N}h_N$, where
$w_ND_N\to J_s^2/\lambda$. No time growing with $N$ is claimed.

\newpage
\section{Six-point consequences and independent finite checks}
The familiar Hamiltonian moment is recovered:
\begin{equation}
 \ph(h^6)=5+6q+3q^2+q^3,\qquad
 \ph(M^3(M^\dagger)^3)=6q^3.
\end{equation}
The actual disorder complement remains
$g=h-\gamma(M+M^\dagger)$, with $\ph(g^2)=1-s$.
Its sixth moment from the joint law is
\begin{equation}\label{six}
 \frac{\ph(g^6)}{(1-s)^3}
 =5+6s+14q+10sq+13q^2+7sq^2+6q^3.
\end{equation}
Together with the preceding fourth moment $2+s+2q$, this
quantifies why a standard independent $q$ mode is not the
Hamiltonian complement.

The new integer CAR driver checks all fifteen three-edge pairings
for $(N,p)=(4,1),(6,1),(4,2)$. There are 1,461,120 independent
Gaussian label assignments. Each raw matrix element is compared
with the single-site product including \eqref{sign}. The first
two systems contain respectively 8 and 72 negative nonzero
contributions; all are reproduced. Independently extracting
\eqref{exact} verifies all $3^3$ general/forward/backward edge
specifications for each pairing and fixture: 1,215 coefficients.
The quadratic local table is checked for all pairings and forward
edge types through four edges, 1,814 cases.

A separate exact symbolic driver lists all 729 sixth-word
polynomials and verifies \eqref{six}. Rational Gram elimination
at $\gamma=1/2,2/3$ gives rank seven through degree two and rank
fifteen through degree three, with positive nonzero pivots.
The latter rank is a two-parameter-point check, not a proof of
generic rank fifteen or a universal binary-tree representation.
Finite row/column covariance counts and the scalar time-tail
majorant are separately recorded. These checks support the
specific implementations; the arbitrary-degree and time claims
rest on the proofs above.

The source, parameters, versions, input hashes and three scientific
outputs are stored in
\texttt{results/susskind-slater-joint-words/}.
The evidence manifest also binds the final LaTeX, PDF, two-pass
build logs and every rendered page. Old successful scientific
outputs were not regenerated or overwritten.

\newpage
\section{What this makes possible next}
The same-$J$ joint dynamics is now specified by a positive
all-word law and a convergent finite-time construction in the
annealed Slater state. It can be used to evaluate dynamical
commutators, spectral measures with selected insertions and
state-preparation kernels. It also supplies a precise object
against which proposed canonical or bulk representations can be
tested, beyond matching two or four moments.

The next constructive task is to turn the Gram/multiplication
description into a transparent transfer representation and to
compute physically useful kernels efficiently. The degree-three
rank check suggests an economical organization, but no such
organization is assumed in this paper. The four second-layer
directions of the preceding stage must all survive any proposed
representation with their known Hamiltonian coefficients.

The microscopic Slater vector is not stationary under $h$.
A physical observer state, its preparation and its relation to
the original static-patch observables must therefore be specified
before identifying particle masses or spatial generators. The
joint moment limit here is annealed; concentration for the full
joint algebra is not asserted. It has fixed positive $\lambda$
and fixed time windows, and says nothing about recurrences at
times growing with $N$. Nor is this a derivation of a finite-$G$
gravitational scattering amplitude. These remaining questions
preserve the original de Sitter research target.

\begin{thebibliography}{9}
\bibitem{BNR} M. Berkooz, V. Narovlansky and H. Raj,
\emph{Complex Sachdev-Ye-Kitaev model in the double scaling limit},
\href{https://arxiv.org/abs/2006.13983v2}{arXiv:2006.13983v2} (2020).
\bibitem{RX} H. Rajgadia and J. Xu,
\emph{Emergent States and Algebras from the Double-Scaling limit
of Pure States in SYK},
\href{https://arxiv.org/abs/2604.14387v3}{arXiv:2604.14387v3} (2026).
\bibitem{Neeb} K.-H. Neeb,
\emph{On Analytic Vectors for Unitary Representations of Infinite
Dimensional Lie Groups},
\href{https://arxiv.org/abs/1002.4792v2}{arXiv:1002.4792v2} (2010),
Theorem~4.8 (Nelson's theorem).
\end{thebibliography}
\end{document}
